\begin{abstract}
Trading card game agents act through variable, multi-stage legal menus under
partial observability. This makes representation correctness inseparable from
policy quality: two options with different public rule semantics may be aliased
by a transient option index, while a richer representation can accidentally
leak an opponent's private state. We describe a direct-policy Pok\'emon TCG
controller that keeps the pinned simulator as the sole legality and transition
authority, restricts inference to the acting player's information set, and adds
a zero-bypass public-rule semantic residual to a preserved causal policy. A
content-addressed 20-day corpus and collision census drive branch activation;
unsupported metadata, auxiliary, and checklist routes remain exact-zero. The
census contains 18,412,973 option records and identifies 127,641 collision
groups in the legacy projection. The repaired key leaves zero unresolved
actionable semantic collisions, with 99 residual duplicate groups classified as
harmless permutations. We further specify full-model bootstrap, direct-policy
self-play, a trainer-only public Prize-plan canary, matched evaluation, and
immutable receipt lineage. A 900-game iteration-2 holdout is reported with its
non-promotion pipeline disposition, and an iteration-3 holdout waiver is
reported as zero measured games. The paper's central claim is therefore
representation and auditability, not leaderboard superiority or causal
gameplay improvement.
\end{abstract}

\begin{IEEEkeywords}
reinforcement learning, partial observability, legal-action representation,
trading card games, reproducibility, semantic collisions, policy audit
\end{IEEEkeywords}
